\documentclass[11pt,a4paper]{article}

\usepackage[margin=1.05in]{geometry}
\usepackage{amsmath,amssymb,amsthm}
\usepackage{mathtools}
\usepackage{stmaryrd}
\usepackage{microtype}
\setlength{\emergencystretch}{2.5em}
\usepackage{enumitem}
\usepackage{booktabs}
\usepackage{array}
\usepackage{xcolor}
\usepackage[colorlinks=true,linkcolor=blue!45!black,urlcolor=blue!60!black]{hyperref}

\newtheorem{theorem}{Theorem}[section]
\newtheorem{proposition}[theorem]{Proposition}
\newtheorem{lemma}[theorem]{Lemma}
\newtheorem{corollary}[theorem]{Corollary}
\newtheorem{conjecture}[theorem]{Conjecture}
\theoremstyle{definition}
\newtheorem{definition}[theorem]{Definition}
\newtheorem{example}[theorem]{Example}
\newtheorem{milestone}{Milestone}
\theoremstyle{remark}
\newtheorem{remark}[theorem]{Remark}
\newtheorem{question}[theorem]{Question}
\newtheorem{claim}[theorem]{Claim}
\newtheorem{observation}[theorem]{Observation}
\newcommand{\HQB}{\texttt{hecke-\allowbreak{}in-\allowbreak{}q-\allowbreak{}bool.pdf}}
\newtheorem{hazard}[theorem]{Hazard}

\newcommand{\FF}{\mathbb{F}}
\newcommand{\NN}{\mathbb{N}}
\newcommand{\ZZ}{\mathbb{Z}}
\newcommand{\CC}{\mathbb{C}}
\newcommand{\PP}{\mathbb{P}}
\newcommand{\OO}{\mathcal{O}}
\newcommand{\cat}[1]{\mathbf{#1}}
\newcommand{\Set}{\cat{Set}}
\newcommand{\Vect}{\cat{Vect}}
\newcommand{\Poly}{\cat{Poly}}
\newcommand{\Cont}{\cat{Cont}}
\newcommand{\Span}{\cat{Span}}
\newcommand{\Prof}{\cat{Prof}}
\newcommand{\Rep}{\mathrm{Rep}}
\newcommand{\ext}[1]{\llbracket #1 \rrbracket}
\newcommand{\con}{\triangleleft}
\newcommand{\dual}[1]{{}^{L}\!#1}
\DeclareMathOperator{\Hom}{Hom}
\DeclareMathOperator{\End}{End}
\DeclareMathOperator{\GL}{GL}
\DeclareMathOperator{\Gr}{Gr}
\DeclareMathOperator{\rank}{rank}
\newcommand{\qbin}[2]{\genfrac{[}{]}{0pt}{}{#1}{#2}_{q}}
\newcommand{\Fl}{\mathrm{Fl}}
\newcommand{\Groth}{\mathfrak{G}}

\definecolor{okgreen}{rgb}{0.0,0.45,0.15}
\definecolor{nored}{rgb}{0.65,0.0,0.0}
\definecolor{openblue}{rgb}{0.0,0.25,0.6}
\newcommand{\GO}{\textcolor{okgreen}{\textbf{do}}}
\newcommand{\STOP}{\textcolor{nored}{\textbf{don't}}}
\newcommand{\WAIT}{\textcolor{openblue}{\textbf{gated}}}

\title{\bfseries Three Directions, and an Objection to the Question\\[0.4em]
\large what to actually do with \HQB{}}
\author{Robin Langer \and Claude}
\date{24 September 2026}

\begin{document}
\maketitle

\begin{abstract}
\HQB\ establishes that \texttt{q-bool} and
\texttt{lascoux} contain, unnamed, a number of standard objects of the Langlands
programme. Every one of those identifications is classical. This document argues
that the obvious continuations --- building the affine flag container, chasing
the dual group through Satake --- are therefore re-derivations rather than
research, and that the one place where the container formalism has a feature the
Hecke literature lacks is the shape/position asymmetry that
\texttt{q-container.pdf} Theorem~B records as an obstruction and that
\S\ref{sec:asym} here reinterprets as a description of a cell decomposition.
Three directions follow from that reading. \textbf{A} (\S\ref{sec:A}):
Hecke correspondences are spans, not maps, so the right container-theoretic home
for them is a bicomodule; RSK wants the same missing structure, which joins the
two repositories at the level of the open problem rather than the theorem.
Highest ceiling, gated on an unobtained paper. \textbf{B} (\S\ref{sec:B}):
expander codes, joining \texttt{error-correcting.pdf} to the Ramanujan-graph face
of the Hecke operator --- tractable, applied, lower ceiling. \textbf{C}
(\S\ref{sec:C}): a paving theorem characterising $q$-containers, which would give
\texttt{q-container.pdf} the central theorem it currently lacks and explain its
own refuted conjecture. \S\ref{sec:not} says what not to do, \S\ref{sec:seq}
gives a sequencing with kill criteria, and \S\ref{sec:priors} gives honest
priors.
\end{abstract}

\section{The objection: recognition is not research}\label{sec:obj}

The mishnah is that \HQB{} opens a programme. Ten links,
four verifications, a list of next steps.

The objection is that all sixteen rows of its status table are marked
\textbf{classical} or \textbf{open-and-not-attempted}, and that the classical
ones are classical by a wide margin: Iwahori 1964, Lascoux--Sch\"utzenberger
1982, Kostant--Kumar 1986, Macdonald's Chapter~V, Bernstein, Cherednik,
Mirkovi\'c--Vilonen. What the document does is \emph{recognition}. It tells you
where you already were. It produces nothing.

The five continuations listed in its \emph{mai nafka minah} inherit that
defect. Building the affine flag container terminates in a theorem of
Bernstein's. Chasing Satake terminates in Macdonald's Chapter~V. Deciding
whether the $t$-adic grading induces Macdonald's $q$ terminates in Cherednik.
Each is a re-derivation in container notation of an object that does not need a
new presentation, and \texttt{lascoux/lenses-in-my-thesis2.pdf} already contains
the correct instinct about this in its box headed \emph{what the container
language does not buy}.

The resolution is to stop asking which link to chase and ask instead: where does
the container formalism have a feature that the Hecke literature lacks? There is
exactly one, and \texttt{q-bool} has already proved a theorem about it.

\section{The asymmetry is a description, not an
obstruction}\label{sec:asym}

\texttt{q-container.pdf} Theorem~B states that linearising the shapes destroys
the theory: a $q$-container must keep its shapes a bare set while its positions
become $\FF_q$-spaces. \texttt{data-control.pdf} elevates this to a principle
and argues, in its \S``Why the distinction cannot be repaired away'', that the
set-ness of shapes is not a limitation awaiting a generalisation.

Both documents treat this as an obstruction discovered. Read against \S3 of
\HQB{}, it is instead a description arrived at.

\begin{observation}\label{obs:asym}
A Bruhat stratification of a flag variety is exactly a \emph{set} of
cells, indexed by the Weyl group $W$, each of which is an \emph{affine
space} $\FF_q^{\ell(w)}$. Shapes bare, positions linear. The reason shapes cannot be
linearised is that $W$ is not a variety: there is no Weyl-group-shaped vector
space for them to become.
\end{observation}

This is worth stating carefully because it changes what the formalism is
\emph{for}. If Observation~\ref{obs:asym} is right in general, then a
$q$-container is not a deformation of a data type. It is a presentation of a
space with an affine paving, and the counting series of
\texttt{q-container.pdf}~\S6 is a Poincar\'e polynomial. All three directions
below are consequences of taking that seriously; \S\ref{sec:C} is the attempt to
prove it.

\section{Direction A: Hecke correspondences as bicomodules}\label{sec:A}

\subsection{The negative result that motivates it}

\HQB{}~\S10 records, among the things that do not link,
that the Hecke correspondence is not a container morphism, because both legs
move shapes. That negative is the interesting part.

A container morphism $(u,f)$ is a \emph{map}: $u$ dispatches forward on shapes,
$f$ routes backward on positions. A Hecke operator is a \emph{span}:
\[
  \Fl_n(\FF_q) \;\xleftarrow{\ p_1\ }\;
  Z_i := \Fl_n \times_{\Fl_n^{(i)}} \Fl_n
  \;\xrightarrow{\ p_2\ }\; \Fl_n(\FF_q),
  \qquad
  Z_i = \{(F,F') : \pi_i F = \pi_i F'\},
\]
and the algebra structure is composition of spans by pullback --- which is
convolution, which is why $T_iT_j$ means anything at all. Spans of $G$-sets are
$G$-equivariant profunctors; the Hecke algebra is an endomorphism algebra in
that bicategory.

So the missing notion is not a better kind of map. It is the layer above maps.

\subsection{What that layer is in this formalism}

Comonoids in $\Poly$ are small categories (Ahman--Uustalu for directed
containers; Garner), and bicomodules between polynomial comonoids are the
profunctor-like layer above them, specialising to spans when the comonoids are
discrete. The formalism therefore already has the right shape of object. What it
does not have is the $\FF_q$-linearised version, which is precisely the
\texttt{q-bool} setting.

\begin{definition}\label{def:qspan}
Provisional. A \emph{$q$-bicomodule} from $S \con P$ to $T \con Q$ consists of a
set $R$ of \emph{correspondence shapes} with maps $\sigma : R \to S$ and
$\tau : R \to T$, together with, for each $r \in R$, a linear relation
\[
  \rho_r \;\subseteq\; P(\sigma r) \oplus Q(\tau r) ,
\]
composed by pullback on shapes and composition of linear relations on positions,
with multiplicities given by counting fibres.
\end{definition}

The move from \emph{map} to \emph{relation} on the position plane is the whole
proposal, and it is forced: a Hecke structure constant is a count of
intermediate objects, which a function cannot record and a relation can.

\subsection{What would count as success}

Three milestones, in order, each falsifiable and the first two cheap.

\begin{milestone}\label{m:A1}
The Hecke span $Z_i$ of \S\ref{sec:A} is a $q$-bicomodule in the sense of
Definition~\ref{def:qspan}, with $R$ the set of pairs $(w, w')$, $w' \in \{w,
ws_i\}$, and $\rho$ the incidence of Schubert cells. One week.
\end{milestone}

\begin{milestone}\label{m:A2}
Composition of $q$-bicomodules reproduces
\[
  T_i^2 \;=\; (q-1)T_i + q\,I .
\]
This is the sharp test, and it is decidable against
\texttt{code/verify.py} V1, which already has the matrices for
$(n,q) \in \{(2,2),(2,3),(3,2),(3,3),(4,2)\}$. If the composition law of
Definition~\ref{def:qspan} does not produce the $q$ in the second term, the
definition is wrong and should be changed rather than defended.
\end{milestone}

\begin{milestone}\label{m:A3}
RSK is a $q$-bicomodule. \texttt{lascoux-container.pdf}~\S``What RSK becomes''
already conjectures that RSK is a bicomodule between the two comonoids, and
lists it as open. If Milestones~\ref{m:A1}--\ref{m:A3} all land, the two
repositories share not a theorem but an open problem, which is a stronger form
of connection: one construction closes both.
\end{milestone}

Beyond that, the ambitious end: whether the monoidal structure on
$q$-bicomodules reproduces the Satake tensor structure, hence
$\Rep(\dual{G})$ as structure on a polynomial comonoid. That is the payoff and
it is not a milestone, because nothing about it is currently plannable.

\subsection{The hazard, which is internal as well as external}

\begin{hazard}\label{haz:licsno}
Composition in Definition~\ref{def:qspan} requires the pullback on the shape
plane to interact coherently with relation-composition on the position plane ---
that is, a distributive law between the two planes of
\texttt{data-control.pdf}. The LICS 2024 no-go theorems on directed-container
distributive laws bear directly on this. They are listed in
\texttt{lascoux/papers/INDEX.md} as \emph{paywalled and not obtained}.
\end{hazard}

This hazard is not merely a missing citation. \texttt{data-control.pdf} argues at
length that the data/control distinction cannot be repaired away; a distributive
law is exactly a repair; so the document's own central principle predicts that
the distributive law may fail. If it does, Direction~A is dead, and the failure
is informative rather than merely disappointing --- it would say that Hecke
convolution is not expressible in this formalism and would explain why.

\textbf{Obtain that paper before committing a month.} It is an afternoon, and it
gates everything above.

\section{Direction B: expander codes}\label{sec:B}

The tractable direction, and the only one with applications.

Three facts already in hand, from three different documents:

\begin{enumerate}[nosep,label=(\roman*)]
\item \texttt{error-correcting.pdf}: a container morphism over $\FF_q$
dispatches on a shape and then routes data \emph{without inspecting it}; the
routing is a matrix, so the morphism is a linear code, composition is matrix
product, and rank is monotone under composition.
\item \HQB{}~\S2: the Hecke operator for $\GL_2$ is the
adjacency operator of the $(q+1)$-regular Bruhat--Tits tree, and Deligne's
bound $|a_p| \le 2p^{(k-1)/2}$ is the optimal spectral gap --- which is where
Ramanujan graphs come from (Lubotzky--Phillips--Sarnak, by arithmetic quotient).
\item Sipser--Spielman expander codes are built from a $d$-regular expander plus
a small inner code, and decode by a \emph{local} flip rule.
\end{enumerate}

\begin{question}\label{q:expander}
Is a Sipser--Spielman code a container morphism whose shape set is the vertex
set of an expander and whose position spaces carry the inner code? If so the
expansion parameter is a property of the \emph{shape} plane and the inner code a
property of the \emph{position} plane, and locality of decoding is the
statement that routing does not inspect.
\end{question}

That last clause is why this direction is worth doing rather than merely
possible. \texttt{data-control.pdf}'s distinction is decorative in most of its
instances; in a locally-decodable code it does real work, because ``the decoder
looks at the graph and not at the data'' is a correctness property, not a
description.

\begin{milestone}\label{m:B1}
Express the Tanner / Sipser--Spielman construction as a composite of container
morphisms, and check that the rank-monotonicity theorem of
\texttt{error-correcting.pdf} yields the standard rate bound
$R \ge 1 - 2(1-r)$ for inner rate $r$. If the formalism cannot recover a bound
it should be able to recover, it is not buying anything here either.
\end{milestone}

\begin{milestone}\label{m:B2}
Exhibit an LPS Ramanujan graph as the shape plane of such a container and read
the spectral gap off the Hecke operator of \HQB{}~\S2
rather than importing it. This is the join: the same operator, in the same
$q$, serving a coding-theoretic purpose.
\end{milestone}

Ceiling: modest. The coding theory is well-trodden and the container reading is
unlikely to improve a bound. Floor: high --- the milestones are achievable and
the write-up is publishable as an exposition even if nothing is new.

\section{Direction C: the paving theorem}\label{sec:C}

Self-contained, improves an existing document rather than starting a new one,
and is the direct attempt to prove Observation~\ref{obs:asym}.

\begin{conjecture}\label{conj:pave}
A $q$-container with finite shape set ``is'' a finite-type $\FF_q$-scheme with
an affine paving indexed by that set: shapes index strata, positions are the
strata, and the counting series of \texttt{q-container.pdf}~\S6 is the
Poincar\'e polynomial of the paving evaluated at the field size.
\end{conjecture}

If this is right it does three things at once for \texttt{q-container.pdf}.

\begin{enumerate}[nosep,label=\arabic*.]
\item It \emph{explains} Theorem~B instead of recording it: shapes cannot
linearise because the index set of a paving is not itself a variety.
\item It explains the refuted Conjecture~C in the same breath --- $\Vect_{\FF_q}$
being self-dual is a statement about the position plane, and the failure of a
second derivative is a statement that the shape plane has no dual.
\item It supplies the Targets-table row for $q$-Pascal for free, since the
Grassmannian container ought to have counting series $\qbin{n}{k}$ by the
Schubert paving of $\Gr(k,n)$, and $q$-Pascal is then the cell recursion rather
than an identity to be verified.
\end{enumerate}

\begin{milestone}\label{m:C1}
Compute the Grassmannian container --- shapes the partitions in a $k \times
(n-k)$ box, positions $\FF_q^{|\lambda|}$ --- confirm the counting series is
$\qbin{n}{k}$, and derive $q$-Pascal from the cell recursion. This is a day, and
it is the smallest possible test of Conjecture~\ref{conj:pave}.
\end{milestone}

\begin{milestone}\label{m:C2}
State the correspondence in both directions and find the first counterexample.
A $q$-container is bare data; a paved scheme has a topology; so the
correspondence is at best an equivalence onto a subcategory, and locating the
boundary is the theorem.
\end{milestone}

Ceiling: modest but real --- a characterisation theorem is a better centre for
\texttt{q-container.pdf} than a list of verified propositions. Probability of
something true and provable: high.

\section{What not to do}\label{sec:not}

\begin{description}[leftmargin=1.6em,style=nextline,nosep]
\item[\STOP{} Build the affine flag container.] This is
\HQB{} Question~5.4 and it is my own suggestion; I am
withdrawing it as a priority. It is the most obvious next step and the least
rewarding: a long construction whose endpoint is Bernstein's description of the
centre, and whose novelty is notational.
\item[\STOP{} Chase Macdonald's second parameter as a construction.] Claim~7.2
of that document is very likely right, and being right costs nothing to record
and everything to re-derive.
\item[\WAIT{} Read Macdonald Chapter~V.] Two days, and it is de-risking rather
than direction: it decides whether \S6 of that document stands on anything.
Worth doing, not worth planning around.
\end{description}

\section{Sequencing, with kill criteria}\label{sec:seq}

\begin{center}
\small
\setlength{\tabcolsep}{4pt}
\begin{tabular}{@{}rlll@{}}
\toprule
& step & cost & kill criterion \\
\midrule
1 & Obtain the LICS 2024 no-go paper & an afternoon & --- \\
2 & Macdonald Ch.~V against \texttt{q-container}~\S9 & two days & --- \\
3 & Direction~C, Milestone~\ref{m:C1} (Grassmannian) & a day
  & counting series $\ne \qbin{n}{k}$ \\
4 & Direction~C, Milestone~\ref{m:C2} (the boundary) & two weeks
  & no clean subcategory after two weeks \\
5 & Direction~A, Milestones~\ref{m:A1}--\ref{m:A2} & two weeks
  & composition does not give $(q-1)T_i + q$ \\
6 & Direction~A, Milestone~\ref{m:A3} (RSK) & a month
  & the distributive law of Hazard~\ref{haz:licsno} fails \\
7 & Direction~B if A dies & --- & --- \\
\bottomrule
\end{tabular}
\end{center}

Step~1 gates steps~5--6 and costs an afternoon, so it goes first regardless of
enthusiasm. Step~3 is deliberately placed before step~5: it is cheap, it is
self-contained, and it is useful to have one thing land before starting the one
that might not.

Direction~B is the fallback rather than the plan, because it survives contact
with a negative result: if Hazard~\ref{haz:licsno} closes Direction~A, nothing
about Direction~B changes.

\section{Priors}\label{sec:priors}

Stated before the work, so they can be scored afterwards.

\begin{center}
\small
\setlength{\tabcolsep}{4pt}
\begin{tabular}{@{}clcc@{}}
\toprule
& claim & prior & if it lands \\
\midrule
\ref{obs:asym} & the shape/position asymmetry is a paving & $70\%$
  & reframes \texttt{q-container} \\
\ref{conj:pave} & $q$-containers are $\FF_q$-paved spaces, on a subcategory
  & $60\%$ & central theorem for \texttt{q-container} \\
\ref{m:C1} & Grassmannian container has series $\qbin{n}{k}$ & $90\%$
  & Targets row, free \\
\ref{m:A1} & the Hecke span is a $q$-bicomodule & $75\%$
  & definition is on the right track \\
\ref{m:A2} & composition gives $(q-1)T_i + q$ & $45\%$
  & the definition is correct \\
\ref{m:A3} & RSK is a $q$-bicomodule & $35\%$
  & the two repositories join \\
--- & Direction~A yields something publishable & $25$--$35\%$
  & a container-native Hecke theory \\
\ref{haz:licsno} & the distributive law is the obstruction & $50\%$
  & informative either way \\
\ref{q:expander} & expander codes are container morphisms & $80\%$
  & applied write-up \\
\ref{m:B1} & the rate bound falls out of rank monotonicity & $55\%$
  & the formalism earns its keep \\
--- & anything here reaches arithmetic Langlands & $<5\%$
  & --- \\
\bottomrule
\end{tabular}
\end{center}

The honest summary of that table: the highest-probability rows are the ones that
improve documents you have already written, and the highest-ceiling row is at
$25$--$35\%$ and gated on a paper you cannot currently read. That is a normal
shape for a research programme and is not an argument against it. It is an
argument for doing step~1 this week.

\section{Sources to obtain}

\begin{description}[leftmargin=1.6em,style=nextline,nosep]
\item[LICS 2024, no-go theorems for directed-container distributive laws.]
Paywalled; listed unobtained in \texttt{lascoux/papers/INDEX.md}. Gates
Direction~A. Highest priority.
\item[Ahman, Uustalu.] \emph{Directed containers as categories}. The comonoid
half of Definition~\ref{def:qspan}.
\item[Garner.] \emph{Polynomial comonads} / the comonoid--category
correspondence, and bicomodules as the layer above.
\item[Spivak, Niu.] \emph{Polynomial Functors: A Mathematical Theory of
Interaction}. Bicomodules in $\Poly$; the ambient calculus for
Definition~\ref{def:qspan}.
\item[Macdonald.] \emph{Symmetric Functions and Hall Polynomials}, Ch.~V. Step~2
of \S\ref{sec:seq}.
\item[Sipser, Spielman.] \emph{Expander codes}, IEEE Trans.\ Inform.\ Theory 42
(1996). Direction~B.
\item[Lubotzky, Phillips, Sarnak.] \emph{Ramanujan graphs}, Combinatorica 8
(1988). Milestone~\ref{m:B2}.
\end{description}

\end{document}
